\documentclass[10pt]{beamer}
\usefonttheme{professionalfonts}
%\usetheme{Rochester}
%\usetheme{AnnArbor}
\usetheme{Szeged}
\useoutertheme{infolines}
\usecolortheme{dolphin}
\usepackage{graphicx}
% \usepackage[english]{babel}
\usepackage{hyperref}
\usepackage{multicol}
\usepackage{color}
%\usepackage{smartdiagram}


%\definecolor{carmine}{rgb}{0.59, 0.0, 0.09}
\usepackage{amsmath} 

\usepackage{graphicx}
\usepackage{tikz}
\usetikzlibrary{shapes.callouts}
\usepackage{multirow}

\usepackage[utf8]{inputenc}

\usepackage{enumerate}
\usepackage[shortlabels]{enumitem}
\usepackage{amsmath,amsthm,amssymb,amsfonts,mathrsfs,latexsym,stmaryrd}
\usepackage{tcolorbox}

\usepackage{color} 
\usepackage{amssymb, amsmath, amsbsy}
\usepackage{adjustbox,lipsum}
\usepackage{placeins}

\usepackage{textpos}
\usepackage{wrapfig}
\usepackage[spanish, es-tabla]{babel}

\usepackage{caption}
\setlength{\parindent}{15pt}
%\usepackage[hyperfootnotes=false]{hyperref}

\usepackage{url}
\usepackage{float}
\usepackage{tikz}
\usetikzlibrary{patterns}


\title[ICS-162]{Econometría ICS-294\\ \small{Clase 17: Variables instrumentales}}

\author{Jocelyn Tapia Stefanoni}
\institute[USM]
{Universidad Técnica Federico Santa María\\ % Your institution for the title page
	\medskip
	\textit{} % Your email address
	
}
\logo{
\includegraphics[width=3cm]{certamenes/logo_usm.png}
}
\date{}
\newcommand\Fontvi{\fontsize{9}{7.2}\selectfont}
\begin{document}
	
\begin{frame}
	\titlepage % Print the title page as the first slide
\end{frame}

\section{Introducción}

\begin{frame}{Introducción}
    \begin{itemize}
        \item En las clases pasadas vimos que cuando omitimos una variable $x_{2}$ que se correlaciona con $x_{1}$ se viola el supuesto $E\left[u \mid x_{1}\right]=0$ y nuestro estimador MCO estara sesgado.
        \item Cuando esto pasa, decimos que $x_{1}$ es endógeno (esta correlacionado con el error).
        \item Controlar: Una solucion es incluir la variable $x_{2}$ en el modelo y controlar por ella, pero en general no contamos con todas las variables.
        \item Entonces, que podemos hacer?
        \item Metodo 1: Metodo de variables instrumentales (IV).
    \end{itemize}
\end{frame}

\section{Endogeneidad y IV}

\begin{frame}{Causas de la Endogeneidad}
    Generalmente observamos problemas de endogeneidad en una de estas 3 situaciones:
    \begin{enumerate}
        \item Variables omitidas: 
        \begin{itemize}
            \item Son variables que están correlacionadas con los $x$ (los regresores) pero no las podemos incluir por ser inobservables
        \end{itemize}
        \item Simultaneidad: $x \rightarrow y$ y $y \rightarrow x$
        \item Error de medida en las $X$
    \end{enumerate}
\end{frame}

\begin{frame}{Endogeneidad}
    El problema de la endogeneidad:
    \begin{figure}
        \centering
        \includegraphics[width=0.3\linewidth]{z SebastianG/Imagenes/Grafico4_1.jpeg}
    \end{figure}
    $$
    E\left[\hat{\beta}_{1}^{M C O}\right]=\beta_{1}+\delta_{1} \beta_{2}
    $$  
\end{frame}

\begin{frame}{Solución 1: Variables Instrumentales}
    Una variable instrumental valida, es una variable $\mathbf{Z}$ que correlaciona con $X$, pero no con el error.
    \begin{figure}
        \centering
        \includegraphics[width=0.5\linewidth]{z SebastianG/Imagenes/Grafico4_2.jpeg}
    \end{figure}
    Formalmente, un instrumento es valido si cumple con DOS supuestos claves:
    \begin{enumerate}
        \item Relevancia: $\operatorname{Cov}(X, Z) \neq 0$. La variable $Z$ afecta $X$. Noten que observamos $Z$ y X por lo que este supuseto es testeable.
        \item Exogeneidad o restriccion de exclusion: $\operatorname{Cov}(Z, U)=0$. Dos interpretaciones:
        \begin{itemize}
            \item El instrumento no correlaciona con el error.
            \item $E[Y \mid X]=E[Y \mid X, Z]$ : El instrumento $Z$ no le pega directamente a $Y$. Solo le pega a Y a traves de su efecto en X .
        \end{itemize}
    \end{enumerate}
\end{frame}

\begin{frame}{Intuición}
    \begin{itemize}
        \item La intuicion es que X tiene dos partes. Una parte exogena y una parte endogena (que esta correlacionada con u).
        \item El instrumento Z, al no estar correlacionado con u, me va a permitir usar solo la parte exogena de X para identificar su efecto en Y .
    \end{itemize}
\end{frame}

\begin{frame}{Ejemplo: el efecto de hacer el servicio militar en el crimen}
    $$
    \textit{crimen}=\alpha+\beta_{1} X+U
    $$  
    $X$: toma valor $1$ si la persona hizo el servicio militar en el pasado.
    \begin{itemize}
        \item $X$ es endógena, $E[U \mid X] \neq 0$. Por que?
        \item Entonces el estimador MCO esta sesgado.
        \item Instrumento $Z$: $1$ si la persona salio sorteada aleatoriamente para hacer el servicio militar. Cumple los supuestos?
        
        \begin{enumerate}
            \item Relevante: $\operatorname{Si}!\  \operatorname{Cov}(Z, X) \neq 0$ porque si salías sorteado, el estado te obligaba a ir.
            \item Exógeno: $\operatorname{Si!}\ \operatorname{Cov}(Z, U)=0$ porque $Z$ depende de un sorteo. No correlaciona con nada, salvo con que vayas al servicio militar.
        \end{enumerate}
        
        \item Entonces, se puede usar Z para identificar el efecto. La idea es que el metodo va a usar la variabilidad de las personas que fueron al militar solo porque fueron (o no) sorteados, y no va a usar la variabilidad de las personas que eligieron ir (endógena).
    \end{itemize}
\end{frame}

\begin{frame}{Buenos y malos instrumentos}
    \begin{itemize}
        \item ¿Cuál de los siguientes $Z$ es un buen instrumento?
        \begin{figure}
            \centering
            \includegraphics[width=0.8\linewidth]{z SebastianG/Imagenes/Grafico4_3.jpeg}
        \end{figure}
        \item  Solo el A es un instrumento valido. Ya que cumple con: 1) Relevancia: $\operatorname{cov}(z, x) \neq 0$ y exogeneidad: 2) $\operatorname{cov}(z, u)=0$.
        \item El caso 2 viola exogeneidad porque el instrumento correlaciona con el error. Tenemos que $\operatorname{Cov}(z, u) \neq 0$. Un ejemplo de esto seria usar como instrumento para servicio militar, $Z=$ inteligencia.
        \item El caso 3 viola relevancia. Si bien el instrumento no correlaciona con el error, tampoco correlaciona con X . Es decir, tenemos que se viola relevancia porque $\operatorname{Cov}(x, z)=0$. Un ejemplo de esto seria usar la hora en la que la persona nació.
    \end{itemize}
\end{frame}

\begin{frame}{Ejemplos de variables endógenas e instrumentos en papers}
    \begin{itemize}
        \item \textbf{Retornos de la educación:}
        $$
        \log \left(\text{salario}\right)_{i}=\beta_{0}+\beta_{1} \text { educ }_{i}+u_{i}.
        $$
        \begin{itemize}
            \item $Z$: dummy trimestre de nacimiento (Angrist and Krueger, 1991)
            \item $Z$: proximidad geográfica con alguna Universidad (Card, 1995)
        \end{itemize}

        \item \textbf{Efecto de instituciones en PBI per cápita:}
        $$
        \operatorname{GDP2015}_{i}=\beta_{0}+\beta_{1} \text { instituciones } 2015_{i}+u_{i}
        $$
        \begin{itemize}
            \item $Z$: mortalidad por enfermedades en el pais cuando llegaron los europeos cientos de años atrás (Acemoglu et al,2001).
        \end{itemize}
    \end{itemize}
\end{frame}

\section{Variables Instrumentales}

\begin{frame}{Implementación del método IV}
    \begin{itemize}
        \item Vamos a ver dos métodos para encontrar el $\hat{\beta}^{I V}$.
        \begin{enumerate}
            \item Estimación por Variables Instrumentales (VI). Hoy.
            \item Mínimos cuadrados en dos etapas (MC2E). Proxima clase.
        \end{enumerate}
    \end{itemize}
\end{frame}

\begin{frame}{Estimación por Variables Instrumentales (VI)}
    \begin{itemize}
        \item Tenemos nuestro modelo de regresión simple:
        $$
        y_{i}=\beta_{0}+\beta_{1} x_{i}+u_{i} \quad \forall i=1 \ldots N
        $$
        \item Problema: $\operatorname{Cov}\left(u_{i}, x_{i}\right) \neq 0$
        \item Tenemos un instrumento $z$ que cumple con:
        \begin{enumerate}
            \item $\operatorname{Cov}\left(u_{i}, z_{i}\right)=0$
            \item $\operatorname{Cov}\left(x_{i}, z_{i}\right)\neq 0$
        \end{enumerate}
        \item Podemos escribir entonces:
        $$
        \begin{aligned}
        \operatorname{Cov}\left(y_{i}, z_{i}\right) & =\operatorname{Cov}(\underbrace{\beta_{0}+\beta_{1} x_{i}+u_{i}}_{y_{i}}, z_{i}) \\
        & =\underbrace{\operatorname{Cov}\left(\beta_{0}, z_{i}\right)}_{0}+\operatorname{Cov}\left(\beta_{1} x_{i}, z_{i}\right)+\underbrace{\operatorname{Cov}\left(u_{i}, z_{i}\right)}_{0} \\
        & =\beta_{1} \operatorname{Cov}\left(x_{i}, z_{i}\right)
        \end{aligned}
        $$
        \item Despejando:
        $$
        \beta_{1}=\frac{\operatorname{Cov}\left(y_{i}, z_{i}\right)}{\operatorname{Cov}\left(x_{i}, z_{i}\right)}
        $$
    \end{itemize}
\end{frame}

\begin{frame}{Relación a nivel poblacional y su análogo en la muestra}
    \begin{itemize}
        \item Ahora que tenemos una relación en la \textbf{población} podemos aplicar un criterio de analogía y llevarlo a la muestra para tener una expresión para el \textbf{estimador} de $\beta$ en la \textbf{muestra}
        \item Usamos las covarianzas muestrales y definimos el $\hat{\beta}_{V I}$
        $$
        \hat{\beta}_{V I}=\frac{S_{y_{i}, z_{i}}}{S_{x_{i}, z_{i}}}=\frac{\sum_{i=1}^{N}\left(y_{i}-\bar{y}\right)\left(z_{i}-\bar{z}\right)}{\sum_{i=1}^{N}\left(x_{i}-\bar{x}\right)\left(z_{i}-\bar{z}\right)}
        $$
        \item Todo esto es observable por lo que podemos estimarlo. (R lo hace por ustedes).
        \item Además, si se cumplen los supuestos del instrumento Z, tenemos que el estimador es un estimador consistente (insesgado) del verdadero $\beta$.
        \begin{align*}
            &\underbrace{\hat{\beta}_{V I} \xrightarrow{p} \beta}\\
            &\text{¡}\hat{\beta}_{V I} \text{ es consistente!}
        \end{align*}
        
    \end{itemize}
\end{frame}

\begin{frame}{Estimación con VI en Modelos de Regresión Múltiple}
    \begin{itemize}
        \item La estimación por IV es fácilmente generalizable a modelos con múltiples regresores
        \item En el curso nos centraremos principalmente en el caso un solo regresor endógeno distintos controles exógenos, aunque existe la posibilidad de tener más de una variable endógena
    \end{itemize}

    \underline{\textbf{Estimación por VI con una endógena y varias exógenas}}
    $$
    \ln \left(\text { salario }_{i}\right)=\beta_{0}+\beta_{1} \text { educ }_{i}+\beta_{2} \text { mujer }+u_{i}
    $$
    \begin{itemize}
        \item Idea: usamos la exógena como su propio instrumento :la matriz de $Z$ incluirá; (i) la constante; (ii) el instrumento para educ; y (iii) mujer
        \item Usamos la fórmula anterior!
        $$
        \hat{\beta}_{V I}=\frac{\sum_{i=1}^{N}\left(y_{i}-\bar{y}\right)\left(z_{i}-\bar{z}\right)}{\sum_{i=1}^{N}\left(x_{i}-\bar{x}\right)\left(z_{i}-\bar{z}\right)}=\left(Z^{T} X\right)^{-1} Z^{T} Y
        $$
        \item donde Z es de nxk, X es de nxk, Y es de nxk. $\beta^{IV}$ de kx1.
    \end{itemize}    
\end{frame}

\begin{frame}{Comparar $\hat{\beta}^{IV}$ con $\hat{\beta}^{MCO}$}  
    $$
    \ln \left(\text {salario}_{i}\right)=\beta_{0}+\beta_{1} \text { educ }_{i}+u_{i}
    $$
    \begin{itemize}
        \item Supongamos que la unica variable omitida es inteligencia.
        \item como inteligencia correlaciona positivamente con educ y salario, uno esperaria
        $$
        E\left[\hat{\beta}^{M C O}\right]>\beta
        $$
        \item Imaginen que tengo un instrumento valido. Entonces,
        $$
        E\left[\hat{\beta}^{\prime V}\right]=\beta
        $$
        \item Entonces, si hago las dos estimaciones, esperaria ver que:
        $$
        \hat{\beta}^{M C O}>\hat{\beta}^{I V}
        $$
        \item porque corregi el sesgo de MCO que, en este caso, era sobreestimar el efecto.
        \item Que pasa si no hay variable omitida, pero educ esta mal medida? Como comparan: $\hat{\beta}^{M C O}$ y $\hat{\beta}^{I V}$ ? $\hat{\beta}^{M C O}<\hat{\beta}^{I V}$ porque elimine el sesgo de atenuacion de MCO
    \end{itemize}
\end{frame}

\begin{frame}{Conclusiones}
    \begin{itemize}
        \item Decimos que una variable es endógena cuando $\operatorname{Cov}(u, x) \neq 0$
        \item En presencia de endogeneidad no se cumple uno de los supuestos de Gauss-Markov, y entonces $\hat{\beta}_{\text {MCO }}$ será:
        \begin{itemize}
            \item Sesgado
            \item Inconsistente
            \item No eficiente
        \end{itemize}
        \item Hay 3 fuentes de endogeneidad frecuentes en economía:
        \begin{enumerate}
            \item Omisión de variables relevantes
            \item Simultaneidad
            \item Errores de Medida
        \end{enumerate}
        \item Para poder estimar consistentemente vamos precisar un \textbf{instrumento}
        \item Un instrumento válido es una tercer variable $z$ que cumple con:
        \begin{itemize}
            \item Exogeneidad: $\operatorname{Cov}(u, z)=0$
            \item Relevancia: $\operatorname{Cov}(x, z) \neq 0$
        \end{itemize}
        \item El metodo de IV nos permite estimar consistentemente $\beta$
    \end{itemize}
\end{frame}

\end{document}